\chapter{DSTRUCT}

This library has as objective to allow easy manipulation of standard
data structures, from basic operations to less standard, while useful
ones.
The library is not architecture optimized, so it is possible that you
can find architecture-dependent libraries that are more efficient for
your own needs, in which case this I would not suggest you to use this
library.
This library is meant to be general-purpose, while complying with the
C language, in order to improve its portability.

\section{Conventions}

Each function of the library has a definition starting with
\verb|ds_|, followed by its main purpose (e.g. \verb|list|,
\verb|sort|,\ldots), '\verb|_|', and is specific function.
If several functions exist for a similar feature, but are specialez
for a specific data type, the function specification will be preceded
by a letter designing the data type.
The main ones are \verb|d| for double, \verb|i| for integer and
\verb|g| for generic (\verb|void| pointer).

\section{Function types}

Several routines rely and standardized definition types.
An important situation is the comparison of two elements with respect to some user-defined ordre.
The comparison functions have to follow to following prototype:
\begin{verbatim}
 typedef int (*DSCompareFunc)(const void *a, const void *b);
\end{verbatim}
By convention, a comparison function should return 0 if \verb|a| is considered equal to \verb|b|
(with respect to the defined order), a strictly positive number if \verb|a| is strictly greater
than \verb|b|, and a strictly negative number if \verb|a| is strictly smaller than \verb|b|.

\section{Linked lists}
\label{sec:list}

A linked \Index{list} is probably one of the most used data structure, so the
library supports the two main variants: singly-linked lists and
doubly-linked lists.
The same functions are defined for these two kinds of list, and the
convention is to prefix the function name by \verb|ds_list| for
a doubly-linked list, and \verb|ds_slist| for a singly-linked list.
Both lists take \verb|void *|, allowing the use of any kind of structure.

\subsection{Singly linked lists}

A singly linked list has the type \verb|DS_SList|.

\subsection{Doubly linked lists}

A singly linked list has the type \verb|DS_List|.

\section{Searching functions}

\begin{verbatim}
int ds_search_double(const void *a, const void *b)
\end{verbatim}
\begin{verbatim}
int ds_search_int(const void *a, const void *b)
\end{verbatim}
\begin{verbatim}
int ds_search_pint(const void *a, const void *b)
\end{verbatim}
\begin{verbatim}
long ds_find_dicho(void *base, int nmemb, size_t size, void *critere,
		   DSCompareFunc search)
\end{verbatim}
\begin{verbatim}
long ds_find_inter(void *base, int nmemb, size_t size, void *critere,
		   DSCompareFunc search)
\end{verbatim}

\section{Sorting and permutation functions}

\subsection{Permutations}\index{permutation}
\label{sec:permutations}

A permutation $p$ is defined as vector of size $n$ whose elements are
integers in $0,\ldots,n-1]$, all being distinct, such that is $a$ is
some array, the application of the permutation produce an array $v$
defined as
\[
  v[i] = a[p[i]].
\]

The function \verb|ds_perm| and \verb|ds_perm_inp| apply some
pertumation to an general array, the first producing a new array, the
last replacing the original array by the permuted one.

\subsection{Heapsort}

We have not implemented the Quicksort algorithm since it is already
available in the standard C library, and in most cases, it is
sufficient (and efficient), but still has a worst case in $O(n^2)$.
We present an implementation of the Heapsort algorithm, both for
standard types and for generic data.
While the implementation is sufficient for most needs, 
C++ programmers are also adviced to turn for the Standard Templates
Library (at least for good implemented ones), since the proposed
functions are often reported to be faster than the C implementations,
including those present in Standard C Libraries.
More efficient, but more difficult to implement, is the Introsort
method, due to Musser~\cite{Muss97}, that combines the usual nice
behaviour of Quicksort, while avoiding the worst case, since its
complexity follows the $O(n \log n)$ of Heapsort.
The \verb|dstruct| does not currenctly propose Introsort.

The basic version of the proposed Heapsort algorithm are the following
ones, dealing with double and integer vectors. In order to indicate
which data type is concerned, the name \verb|heapsort| is simply
prefixed by '\verb|d|' (for double) and '\verb|i|' (for integer).
\begin{verbatim}
void ds_sort_dheapsort(int n, double *array);
\end{verbatim}
\begin{verbatim}
void ds_sort_iheapsort(int n, int *array);
\end{verbatim}
Each one takes as arguments the number of elements in the array to
sort, and a pointer to the array itself.

The generic version of Heapsort in \verb|dstruct| has a similar
declaration than the quicksort standard library, so you can easily
change the called function:
\begin{verbatim}
void ds_sort_gheapsort(void *base, size_t nmemb, size_t size,
		       int(*compar)(const void *, const void *));
\end{verbatim}

Sorting is nothing else that generating a particular permutation of
the original vector, and it sometimes useful to knows this
permutation.
Each Heapsort algorithm has an equivalent version generating this
permutation, defined as in Section~\ref{sec:permutations}.

\begin{verbatim}
void ds_sort_dheapsort_with_perm(int n, double *array, int *perm);
\end{verbatim}

\begin{verbatim}
void ds_sort_iheapsort_with_perm(int n, double *array, int *perm);
\end{verbatim}
